\chapter{文献综述与理论基础}

\section{AIGC 与 AI 文旅短视频}

AIGC 是指基于生成式人工智能技术自动或半自动生成文本、图像、音频、视频等内容的技术与应用形态。与传统数字内容生产相比，AIGC 的突出特征在于能够根据提示词、数据输入或语义指令生成新的内容，并支持跨模态组合\cite{lit01}。旅游行业人工智能研究表明，AI 在旅游企业中已广泛用于服务自动化、个性化推荐、数据分析和客户互动，这为理解 AI 文旅短视频的产业背景提供了基础\cite{lit24}。

从文旅传播角度看，AIGC 的应用贯穿声音、画面和整体叙事等多个层面。声音层面的应用包括 AI 配音、虚拟主播解说和合成音频；画面层面的应用包括 AI 生成景观、虚拟场景重构、历史人物复原和城市视觉再创作；整体叙事层面的应用则体现在借助 AI 将地方文化符号、旅游资源和平台流行表达组合起来，形成具备娱乐性或沉浸感的内容。相关研究表明，AI 驱动叙事能够通过情境再现和故事化表达增强用户对文化旅游场景的参与感\cite{lit19}。

进一步来看，AIGC 对文旅短视频的影响并不仅仅体现在“是否生成了一段画面”，还体现在内容生产流程的整体重构上。传统文旅短视频往往需要经历选题策划、脚本撰写、实地拍摄、后期剪辑、配音配乐和平台发布等多个环节，而 AIGC 的介入使得部分环节出现明显压缩。例如，基于提示词的图像生成与视频生成可以快速完成前期视觉方案验证，AI 配音和数字人技术则可以降低真人出镜与录音成本，智能剪辑工具还能够在后期阶段自动完成镜头拼接、字幕生成和节奏匹配。对人手有限、更新频率要求较高的地方文旅账号而言，这种流程层面的变化具有较强现实吸引力。

同时，AIGC 也改变了文旅内容的表达边界。过去许多文旅宣传依赖真实场景拍摄，因此较难在有限预算和拍摄条件下呈现高度奇观化、故事化或时空重构式的画面。AIGC 提供了一种新的可能，即在保留地方景观和文化符号基础上，进一步叠加想象性视觉包装，将“真实景点展示”转化为“地方意象再叙事”。这使文旅短视频不再只是静态景观的线性介绍，而更容易演化为带有剧情感、角色感和平台梗属性的传播内容。也正因为如此，AI 文旅短视频具备了比一般旅游介绍视频更强的形式创新空间。

但需要指出的是，AIGC 并不天然等于高质量传播。技术门槛降低虽然提升了内容生产速度，却也可能带来风格同质化、画面失真、叙事模板化和地方特色弱化等问题。尤其是当多个账号同时使用类似模型、类似提示词或类似流行模板时，平台中的 AI 内容可能迅速出现“看起来很新鲜，但彼此差异不大”的情况。这意味着研究 AI 文旅短视频时，不能只关注其技术先进性，更要关注其最终是否形成差异化的地方表达和可感知的内容质量。

本文所称“AI 文旅短视频”，是指文旅官方账号或地方文旅传播主体在短视频平台发布的相关视频。其内容通常围绕旅游目的地、城市形象、文化资源或旅游体验展开，并在声音、画面或整体视听呈现中使用 AI 生成或 AI 辅助生成技术。该概念强调两个条件：一是内容主题具有文旅属性，二是视频生产或呈现过程中存在可识别的 AI 介入。

在已有 AIGC 研究中，内容质量是影响用户接受和传播的重要因素。AI 生成视频的视觉和谐性、文本一致性和整体质量会影响用户对内容的理解与评价\cite{lit03}。因此，本文将“AI 技术质感”作为关键变量之一，用于衡量 AI 内容最终呈现的完成度和视觉质量。该变量不同于“AI 介入程度”。介入程度反映生产过程中的 AI 参与深度，而技术质感反映用户最终看到的内容质量。

从产业实践层面看，AI 文旅短视频之所以迅速出现，还与地方文旅部门的传播任务变化有关。近年来，各地文旅账号不仅承担常规资讯发布功能，还需要在节假日营销、热点联动、城市形象更新和年轻受众沟通中持续生产更具话题性和视觉冲击力的内容。在这种高频更新压力下，AIGC 工具为地方文旅传播提供了一种“更快完成视觉包装”和“更容易制造差异化表达”的路径。也就是说，AIGC 并不是孤立的技术潮流，而是在地方文旅传播竞争加剧、平台内容卷入度提高的背景下被迅速吸纳进运营实践的。

不过，当前 AIGC 在文旅内容中的应用也表现出明显的不均衡性。一类账号更多把 AI 用作局部辅助工具，例如 AI 配音、AI 海报转视频、AI 镜头补帧和 AI 标题生成；另一类账号则直接把 AI 作为主要视觉生产方式，尝试生成超现实城市景观、拟人化地标建筑和重构式文化叙事。不同使用方式对应的传播效果未必一致。前者可能更贴近日常运营，后者则更容易形成话题和视觉奇观。也正因如此，本文不将 AI 内容视为单一同质对象，而是希望通过变量区分不同形态的 AI 介入特征。

\section{文旅短视频传播研究}

短视频平台改变了旅游信息传播方式。相较于传统图文和长视频，短视频具有时长短、节奏快、视觉冲击强、互动门槛低和算法推荐明显等特点，更容易在用户碎片化时间中完成目的地展示和情绪唤起。既有研究表明，旅游短视频内容营销会通过感知价值、情感反应和目的地态度影响用户旅游意图\cite{lit04}；TikTok 等平台上的短视频因素也会影响不同代际用户的旅游行为意图\cite{lit05}。

文旅短视频传播效果受到多种内容因素影响。部分研究强调视觉内容、叙事吸引力和情绪唤起在短视频营销中的作用，认为更具吸引力的视觉表达能够提升用户对目的地的兴趣\cite{lit06}。也有研究指出，短视频内容会通过信息性、娱乐性和感知价值影响用户旅游意图\cite{lit07}。从内容类型出发的研究则说明，不同视频风格和表达方式会通过“受启发”和“想行动”的连续心理过程影响潜在游客的出行意图\cite{lit22}。

从平台传播机制看，文旅短视频并不是在真空中传播的。短视频平台的推荐逻辑强调前几秒吸引力、完播倾向、互动反馈和后续扩散潜力，这使得文旅内容必须在“信息价值”与“平台适配性”之间取得平衡。对传统文旅宣传而言，完整介绍景点历史、区位、交通和游览攻略可能有较高的信息密度，但在短视频平台中，如果缺乏足够鲜明的视听钩子，内容就可能在很短时间内被用户划走。因此，研究文旅短视频时，不能仅从旅游信息传播角度理解其效果，也必须把平台算法环境和用户即时注意力竞争纳入分析。

在地方形象塑造层面，文旅短视频还承担着“把地方讲给外部受众听”的功能。城市地标、历史文化、地方方言、美食风味和生活氛围，都可以在短视频中被压缩为一组可快速识别的符号。对于尚未到访该地的用户而言，短视频往往构成其形成第一印象的媒介界面。也就是说，文旅短视频并不仅仅是景点介绍工具，它同时参与地方想象、目的地品牌和城市情感标签的建构过程。AI 技术的介入会进一步改变这种建构方式，因为它能够放大、重组甚至重新创造地方符号。

地方传播视角下的研究进一步强调，短视频嵌入文旅融合之后，不仅承担旅游信息发布功能，也成为地方文化符号再生产和城市形象塑造的重要工具\cite{lit09}。以城市为对象的研究显示，文旅短视频能够影响城市形象评价，并通过内容呈现和用户感知影响传播效果\cite{lit10}。

不过，现有文旅短视频研究大多聚焦一般短视频内容，对 AI 生成或 AI 辅助生成内容的细分特征关注不足。AI 文旅短视频既继承了文旅短视频的视觉传播逻辑，又具有生成式内容的新奇性、非现实性和技术痕迹。因此，仅依靠传统变量难以完整解释传播差异。本文在既有研究基础上，引入 AI 参与方式、呈现质量和视频风格等变量，以更贴近用户实际感知的方式分析传播效果。

此外，现有研究在样本选择上也存在一定局限。许多研究采用热门案例、单个城市样本或实验式内容刺激，虽然有助于进行深描，但对于“官方文旅账号在真实平台中如何批量生产和发布 AI 内容”这一问题的回答仍不充分。本文选取省级文旅账号样本，正是希望在保持样本主体相对一致的基础上，观察不同地区、不同内容策略和不同 AI 呈现方式之间的传播差异。这一处理使研究更接近真实运营场景，也为后续的实证分析提供了较强现实基础。

还需要指出，文旅短视频传播研究中一个常见问题是容易把内容效果过度归因于旅游资源本身。例如，拥有知名山水、强地标和高热度城市形象的地区，往往天然更容易获得关注；但在短视频平台中，资源优势并不会自动转化为传播优势。同样的山水、同样的古建、同样的美食，在不同表达方式下可能产生完全不同的互动结果。AI 技术的出现进一步放大了这种差异，因为它既能够帮助地方资源“被看见”，也可能由于模板化表达导致资源特色被稀释。本文后续把视频风格、内容导向和技术质感放在同一框架中考察，正是为了避免把传播效果简单归因为地方资源禀赋。

\section{社交媒体用户参与与传播效果}

在社交媒体营销研究中，用户参与是衡量传播效果的重要指标。用户参与不仅包括心理层面的关注、兴趣和情感投入，也包括行为层面的点赞、评论、分享、收藏等可观察互动\cite{lit11}。相关研究强调，这些互动行为能够反映受众对内容和品牌的响应程度\cite{lit12}。

在短视频平台中，点赞、评论、收藏和转发具有不同含义。点赞通常代表用户的即时认可，评论反映表达和讨论意愿，收藏体现内容的保存价值，转发则体现其扩散价值。TikTok 短视频推荐研究表明，用户互动行为与平台内容推荐和观看体验密切相关\cite{lit21}。在缺乏后台曝光量、完播率和点击率等数据的情况下，公开可见的互动指标是衡量短视频传播效果的重要替代指标。

本文将点赞量、评论量、收藏量和转发量相加，构造“总互动量”作为传播效果的综合指标。这一做法能够同时考虑用户认可、表达、保存和扩散行为。考虑到互动量通常呈现明显右偏分布，少数爆款视频互动极高，本文在主模型中使用 $\log(\text{总互动量}+1)$，以降低极端值对模型估计的影响。

从传播含义看，这四类互动行为并不完全等价。点赞的门槛最低，更接近用户对内容的即时肯定；评论需要更高的表达投入，因而在一定程度上反映内容引发讨论的能力；收藏通常与“以后还想再看”“对我有参考价值”相关；转发则意味着用户愿意将内容带入新的社交关系链条。将四者综合为总互动量，虽然会损失部分行为差异信息，但能够较好地表征内容在平台生态中的总体传播表现。与此同时，本文在稳健性检验中分别替换不同互动指标作为因变量，也是为了避免总互动量这一综合指标掩盖不同行为之间的细微差异。

还需要看到，用户参与不仅是内容结果，也会反过来影响内容的后续分发。短视频平台通常会根据初期互动信号判断内容的潜在质量，并据此决定是否继续推荐。这意味着点赞、评论、收藏和转发既是传播效果的结果，也是平台推荐链路中的中间反馈变量。对文旅账号而言，提高互动量不仅意味着内容更受欢迎，也可能意味着获得更多后续曝光。正因为公开后台数据难以获取，互动指标在现实研究中才更具有替代价值。

需要说明的是，平台互动效果并不等同于真实旅游转化。用户点赞或转发视频，并不必然意味着实际出游或消费。但对本文关注的“短视频营销效果”而言，互动量能够反映内容在平台传播中的吸引力和扩散能力，是可以观察、可以量化且与研究问题直接相关的指标。

从研究操作层面看，以公开互动指标作为结果变量还有一个现实优势，即可比性较强。不同省份文旅账号的后台统计标准、投流策略和内部运营口径难以统一，而点赞、评论、收藏和转发作为平台前台可见数据，至少在同一平台内具备较稳定的可观测性。虽然这种做法无法完全替代真实转化指标，但对于基于公开平台数据展开的研究而言，它既能够保证样本覆盖面，又能够支撑后续差异检验和回归分析，具有较好的可实施性。

\section{AI 真实感、信任与用户接受}

AI 生成内容虽然能够提升内容生产效率和视觉表现力，但用户对 AI 内容的真实感、可信度和情感接受度仍然会影响互动行为。相关研究表明，用户会依据内容来源、表达质量和自身认知判断 AI 内容是否可信\cite{lit14}。在旅游场景中，这一问题尤为重要，因为旅游内容会影响用户对目的地的想象和信任，如果 AI 内容过度脱离真实场景，就可能削弱目的地可信度。

AIGC 营销内容对用户参与的影响并非单纯来自技术本身，而是与价值感知、情感共鸣和用户信任有关\cite{lit18}。在文化旅游叙事中，AI 可以通过故事化表达和沉浸式场景增强用户参与，但这种增强需要建立在内容与地方文化、目的地特色之间的有效连接上\cite{lit19}。如果 AI 只是形式包装，缺乏真实在地性或文化关联，就可能难以转化为稳定互动。

在文旅语境中，真实感并不一定等于百分之百写实，而更接近“用户是否愿意相信这段内容能够代表某种地方印象”。例如，超现实风格的 AI 视频虽然并不追求摄影式写实，但如果其视觉符号、色彩意象和文化元素仍然能够被用户识别为“这是某座城市”“这是某个景区的气质”，它仍然可能形成较高的地方认同和传播吸引力。相反，如果 AI 内容虽然技术上复杂，却与地方文化脱节，用户就容易把它当作泛化的视觉特效，而不是有明确目的地指向的文旅内容。由此可见，文旅内容中的真实感本质上是一种“文化和场景上的可接受性”。

信任问题也值得特别强调。官方文旅账号与普通娱乐账号不同，其内容不仅承担吸引关注的任务，也在一定程度上代表地方公共传播形象。若 AI 生成内容出现明显失真、过度夸张甚至与地方现实严重不符，虽然短期内可能带来一定新鲜感，但也可能损害账号专业性和目的地可信度。尤其是在旅游决策与实际出行相关的场景中，用户既希望看到有创意的表达，也希望内容与真实体验之间保持基本一致。正因如此，本文把“AI 技术质感”作为核心变量之一，实质上也是在考察 AI 内容是否能够同时满足新奇感与可信度两方面要求。

在线旅游规划研究表明，生成式 AI 在旅游规划中的说服效果会受到信息呈现方式和用户判断的影响\cite{lit27}。关于 AI 中介旅游发现的研究也指出，用户参与深度与预订意向之间并非简单线性关系，互动过程中的信任和价值感知具有重要作用\cite{lit28}。这些研究为解释“AI 介入程度未必显著”提供了理论支撑。深度 AI 介入并不必然带来更高传播效果，因为用户最终评价的是内容是否好看、可信、有趣和值得互动，而不是生产过程中 AI 用了多少。

因此，本文没有将 AI 使用简单处理为一个二元变量，而是把技术使用过程和用户可感知的呈现结果分开考察。声音、画面或混合生成反映 AI 介入的载体，浅度到深度的差异反映 AI 参与生产的程度，画面完成度和视觉自然度则体现用户最终接触到的技术质感。相比生产过程本身，用户更容易依据画面质量、风格新颖性和真实感做出互动反应。

从更广的视角看，AI 真实感与信任问题也连接着地方形象治理。对于官方文旅账号而言，内容不仅面向流量竞争，也面向公共传播责任。若内容虽然短期出圈，却因失真、夸张或与地方认知严重不符而引发质疑，那么这种传播未必真正有利于目的地形象积累。由此可见，AI 内容在文旅场景中的有效传播，并不是“越炫越好”或“越像特效越好”，而是要在奇观化表达、文化辨识度与地方可信度之间保持平衡。本文的研究逻辑，本质上也是试图从传播结果反推这种平衡大致体现在哪些可观察内容变量上。

\section{理论基础：S-O-R 模型}

S-O-R 模型由“刺激（Stimulus）-机体（Organism）-反应（Response）”三个环节构成，常用于解释外部环境刺激如何通过个体心理感知影响行为反应。旅游短视频内容营销研究常借助这一思路分析内容刺激、用户感知与行为意图之间的关系，强调内容特征并非直接产生行为结果，而是需要经过用户感知和情绪反应的中介过程\cite{lit04}。对于短视频传播而言，视频内容、画面风格、声音形式和平台呈现可以视为外部刺激；用户的审美愉悦、新奇感、真实感、信任感和情绪唤起可以视为内部心理状态；点赞、评论、收藏、转发等互动行为则是外显反应。

本文借助 S-O-R 模型构建 AI 文旅短视频传播效果分析框架。在刺激层面，本文关注 AI 介入形式、AI 介入程度、AI 技术质感、视频主体、视频主题、内容导向、视频风格、视频时长和发布距采集日天数。其中，AI 技术质感和视频风格直接影响用户对视频的视觉评价；视频主题和内容导向影响用户对内容意义和信息价值的理解；视频时长影响观看成本；发布距采集日天数影响互动积累时间。

在机体层面，本文并未直接通过问卷测量用户心理变量，而是结合既有文献进行理论解释。可以推断，精良 AI 质感可能提升用户审美愉悦和可信度，抽象/超现实风格可能增强新奇感和分享意愿，过长视频则可能增加观看成本并降低互动概率。反应层面则对应本文的传播效果指标，即点赞、评论、收藏、转发和总互动量。

\begin{figure}[!htbp]
  \centering
  \includegraphics[width=\textwidth]{../logs/extended_work/05_SOR理论分析框架.png}
  \caption{AI 文旅短视频营销效果的 S-O-R 分析框架}
  \label{fig:sor}
\end{figure}

基于该框架，本文并不将 AI 技术视为自动带来传播优势的因素，而是强调 AI 技术需要通过用户可感知的内容质量、视觉风格和观看体验发挥作用。后续研究设计和实证分析均围绕这一逻辑展开。

\section{本章小结}

本章围绕 AIGC、文旅短视频传播、社交媒体用户参与、AI 真实感与信任以及 S-O-R 模型进行了梳理。综合既有研究可以看出，AI 文旅短视频的传播效果既受技术介入方式影响，也受最终内容质量、视觉风格、地方文化连接和平台注意力机制共同制约。现有研究虽然为理解文旅短视频传播和 AIGC 内容接受提供了重要基础，但针对“官方文旅账号”“真实平台样本”“AI 介入特征细分”三者交叉场景的实证研究仍相对有限。

基于此，本文将后续研究重点放在两个层面：一是从内容刺激层面区分 AI 介入形式、AI 介入程度、AI 技术质感以及视频主题、内容导向和风格等变量；二是从传播结果层面以公开互动指标表征用户反应。S-O-R 模型为这一分析提供了理论连接，使本文能够在不直接测量用户心理变量的情况下，从外部内容特征、潜在用户感知和最终互动行为之间的逻辑关系出发构建研究设计。下一章将在此基础上进一步说明研究假设、样本筛选、变量编码和模型设定。
